\section{Conclusion}

\subsection{Bilan intermédiaire}

À mi-parcours du projet, l'équipe a figé l'architecture globale du
scanner (mécanique, électronique, logiciel), validé les choix
techniques principaux et produit une version fonctionnelle du
pipeline logiciel en simulation. Les objectifs SMART définis en
introduction restent atteignables dans le cadre du budget
(\textbf{E16}) et du délai de livraison du 7\,juin\,2026
(\textbf{E22}).

\subsection{Revue des objectifs}

\begin{table}[H]
  \centering
  \begin{tabular}{@{}p{1cm}p{9cm}p{4cm}@{}}
    \toprule
    \textbf{Réf.} & \textbf{Objectif} & \textbf{Statut} \\
    \midrule
    O1  & Numériser un objet $\leq 150^3$\,mm, automatisé          & En cours (intégration) \\
    O2  & Export STL ou OBJ                                         & Fait (logiciel) \\
    O3  & Cycle de scan $\leq 10$\,min                              & À mesurer \\
    O4  & Prise en main $\leq 10$\,min, autonomie $\geq 90$\,\%     & À valider (essais utilisateurs) \\
    O5  & Budget $\leq 300$\,CHF                                    & Piloté \\
    O6  & Sécurité utilisateur (laser, mobile, matériaux)           & Conçu, à valider sur prototype \\
    O7  & Retour visuel LED / écran                                 & Fait (logiciel) \\
    O8  & Interface USB ou web, scan en $\leq 5$ actions            & Fait (web Flask + SSE) \\
    O9  & Erreur $\leq 2$\,mm, répétabilité $\leq 1$\,mm            & À mesurer (risque principal) \\
    O10 & Encombrement $\leq 500^3$\,mm, masse $\leq 15$\,kg        & En cours (mécanique) \\
    O11 & Objets sans cavité fermée, contre-dépouille $\leq 45^\circ$ & Accepté (limite triangulation mono-vue) \\
    \bottomrule
  \end{tabular}
  \caption{Revue des objectifs au rapport intermédiaire.}
  \label{tab:revue-obj}
\end{table}

\subsection{Résultats principaux}

\begin{itemize}
  \item Architecture logicielle modulaire testée et validée en
        simulation ; pipeline \emph{Acquisition $\rightarrow$
        Traitement $\rightarrow$ Reconstruction $\rightarrow$ Export}
        opérationnel.
  \item Schéma électrique figé, sécurité laser intégrée par
        construction (coupure en état \texttt{ERROR}, interlock
        porte).
  \item Répartition claire des responsabilités dans l'équipe et
        planification cohérente avec la date de livraison.
\end{itemize}

\subsection{Risques identifiés}

\begin{itemize}
  \item \textbf{Précision E24 (O9)} : risque technique principal.
        Mitigation : calibration soignée, moyennage temporel,
        filtrage KDTree, caractérisation du banc de mesure.
  \item \textbf{Intégration mécanique tardive} : impact sur les
        mesures de précision et l'essai utilisateur.
        Mitigation : produire un prototype mécanique simplifié au
        plus tôt pour débloquer l'intégration.
  \item \textbf{Approvisionnement des composants} : à surveiller
        pour ne pas retarder la phase T30.
\end{itemize}

\subsection{Perspectives}

Phase T30 (semaines 9 à 14) :

\begin{enumerate}
  \item Assemblage mécanique et câblage du prototype.
  \item Calibration caméra et plan laser sur matériel réel.
  \item Premier scan bout-en-bout, confrontation aux objectifs
        O3 (temps) et O9 (précision).
  \item Essai d'ergonomie (O4) et ajustements d'interface.
  \item Mesures de consommation et validation thermique.
\end{enumerate}

Phase T40 (semaines 13 à 16) : préparation des livrables (rapport
final, vidéo de démonstration, présentation).

\subsection{Commentaire personnel}

\emph{À compléter en fin de projet (retour d'expérience de
l'équipe).}
